\emph{证明系统}提供纯句法的方法，用以刻画语句之间的后承与相容性。
\emph{相继式演算}就是这样一种证明系统。
其中的一个\emph{推导}由一棵相继式树构成（相继式 $\Gamma \Sequent \Delta$
由~$\Sequent$ 分隔的两个公式序列组成）。
推导中最顶端的相继式是形如 $!A \Sequent !A$ 的\emph{初始相继式}。
要使推导正确，所有其他相继式都必须由某一条\emph{推理规则}正确证成。
这些规则成对出现；对于每个联结词和量词，分别有一条规则作用于相继式的
左侧和右侧。
例如，如果相继式 $\Gamma \Sequent \Delta, !A \lif !B$ 由
$\RightR{\lif}$ 规则证成，
则其上方的相继式（即\emph{前提}）必须是 $!A, \Gamma \Sequent \Delta, !B$。
有些规则还允许调整相继式中语句的顺序或数量，
例如，$\RightR{\Exchange}$ 规则允许交换相继式右侧的两个公式。

如果存在相继式 $\quad \Sequent !A$ 的推导，则称 $!A$ 是一个\emph{定理}，
记作~$\Proves !A$。
如果存在 $\Gamma_0 \Sequent !A$ 的推导，且 $\Gamma_0$ 中的每个 $!B$
都属于~$\Gamma$，
则称 $!A$ \emph{从~$\Gamma$ 可推导}，记作 $\Gamma \Proves !A$。
如果存在 $\Gamma_0 \Sequent \quad$ 的推导，且 $\Gamma_0$ 中的每个 $!B$
都属于~$\Gamma$，
则称 $\Gamma$ \emph{不一致}，否则称其\emph{一致}。
这些概念相互关联，例如，$\Gamma \Proves !A$ 当且仅当
$\Gamma \cup \{\lnot !A\}$ 不一致。
它们也与相应的语义概念相关，例如，若 $\Gamma \Proves !A$，
则 $\Gamma \Entails !A$。
证明系统的这一性质——即从 $\Gamma$ 可推导出的都必然是 $\Gamma$
的语义后承——称为\emph{可靠性}。
\emph{可靠性定理}通过对推导的长度进行归纳来证明，即证明每一步推理
只要其前提相继式有效，就能保持结论相继式的有效性。